\documentclass[11pt]{article}
\usepackage{inote}
\inotesetup{mode=note, qr=true,
  brand={Interactive Note}, subject={HSC Physics 1st Paper · Chapter-02},
  chapter={ভেক্টর}, chapter-en={Vector}, chapter-no={দ্বিতীয় অধ্যায়},
  video=SAMPLEVIDEO}
\begin{document}
\chapterbanner

\sectionbar{Lecture Map — টপিক অনুযায়ী ভাগ}
\begin{tblr}{colspec={X[2.2,l,m] Q[c,m] Q[c,m] X[2.6,l,m] Q[c,m]},
  hlines={gridgray,.5pt},vlines={gridgray,.5pt},
  row{1}={bg=headgray,font=\labelfont\bfseries,halign=c},rowsep=2.5pt}
টাইপের নাম & গুরুত্ব & ভিডিওতে & যা যা আছে & পৃষ্ঠা\\
\textbf{Type-01:} ভেক্টরের লব্ধি & \stars{5} & 03:10 & সামান্তরিক সূত্র, লব্ধির দিক, ২টি ম্যাথ & ২\\
\textbf{Type-02:} আপেক্ষিক বেগ (নদী–নৌকা) & \stars{4} & 18:40 & সর্বনিম্ন পথ ও সর্বনিম্ন সময়, ১টি সৃজনশীল & ৩\\
\textbf{Type-03:} ডট ও ক্রস গুণন & \stars{5} & 31:05 & কোণ নির্ণয়, লম্ব হওয়ার শর্ত, ১টি ম্যাথ & ৩\\
\end{tblr}

\fancyheading{গুরুত্বপূর্ণ সূত্রাবলি ও বিশ্লেষণ}
\begin{formulagrid}[0.52]
\sqitem{সামান্তরিক সূত্র (লব্ধির মান):}
$R=\sqrt{P^2+Q^2+2PQ\cos\theta}$\par
\hspace*{3mm}$\theta=$ ভেক্টরদ্বয়ের মধ্যবর্তী কোণ\par
\sqitem{লব্ধির দিক ($\vec P$ এর সাথে):}
$\tan\alpha=\dfrac{Q\sin\theta}{P+Q\cos\theta}$
&
\begin{tikzpicture}[scale=.78]
  \coordinate (O) at (0,0); \coordinate (P) at (3.4,0);
  \coordinate (Q) at (1.2,1.9); \coordinate (R) at (4.6,1.9);
  \draw[guide] (P)--(R) (Q)--(R);
  \draw[vec] (O)--(P) node[midway,below,lbl]{$\vec P$};
  \draw[vec] (O)--(Q) node[midway,left,lbl]{$\vec Q$};
  \draw[vec,accent] (O)--(R) node[pos=.7,above left,lbl]{$\vec R$};
  \pic[draw,angle radius=6mm,"$\theta$",angle eccentricity=1.45]{angle=P--O--Q};
  \pic[draw,accent,angle radius=12mm,"$\alpha$",angle eccentricity=1.2]{angle=P--O--R};
\end{tikzpicture}\\
\sqitem{উপাংশ:}
$A_x=A\cos\theta,\quad A_y=A\sin\theta$\par
$A=\sqrt{A_x^2+A_y^2},\quad \tan\theta=\dfrac{A_y}{A_x}$
&
\begin{tikzpicture}[scale=.8]
  \draw[vecthin] (-.2,0)--(3.4,0) node[right,lbl]{$x$};
  \draw[vecthin] (0,-.2)--(0,2.1) node[above,lbl]{$y$};
  \draw[vec,accent] (0,0)--(2.8,1.7) node[above right,lbl]{$\vec A$};
  \draw[guide] (2.8,0)--(2.8,1.7) (0,1.7)--(2.8,1.7);
  \draw[vec] (0,0)--(2.8,0) node[midway,below=1pt,lbl]{$A_x$};
  \draw[vec] (0,0)--(0,1.7) node[midway,left,lbl]{$A_y$};
  \draw (.9,0) arc (0:31.3:.9) node[pos=.55,right,lbl]{$\theta$};
\end{tikzpicture}\\
\sqitem{ডট গুণন (স্কেলার):}
$\vec A\cdot\vec B=AB\cos\theta=A_xB_x+A_yB_y+A_zB_z$\par
\sqitem{ক্রস গুণন (ভেক্টর):}
$|\vec A\times\vec B|=AB\sin\theta$
&
\textbullet\ $\vec A\cdot\vec B=0$ হলে ভেক্টরদ্বয় পরস্পর \textbf{লম্ব}।\par
\textbullet\ $\vec A\times\vec B=\vec 0$ হলে ভেক্টরদ্বয় \textbf{সমান্তরাল}।\par
\textbullet\ $\cos\theta=\dfrac{\vec A\cdot\vec B}{AB}$ — কোণ বের করার সবচেয়ে সহজ উপায়।\\
\end{formulagrid}

\begin{unitbox}[দ্রুত মনে রাখো]
\item $\theta=0^\circ$ হলে $R_{\max}=P+Q$;\quad $\theta=180^\circ$ হলে $R_{\min}=|P-Q|$
\item $\theta=90^\circ$ হলে $R=\sqrt{P^2+Q^2}$
\item $\hat i\cdot\hat i=1,\ \hat i\cdot\hat j=0;\quad \hat i\times\hat j=\hat k,\ \hat j\times\hat k=\hat i,\ \hat k\times\hat i=\hat j$
\end{unitbox}

\begin{multicols}{2}
% ---------------------------------------------------------------- TYPE 01
\begin{typebox}{01}{ভেক্টরের লব্ধি}
\withqr{190}{\textbf{কনসেপ্ট:} দুটি ভেক্টর একই বিন্দুতে কোনো কোণে ক্রিয়া করলে তাদের সম্মিলিত ফলকে \textbf{লব্ধি} বলে। দুটি ভেক্টরকে সামান্তরিকের দুই সন্নিহিত বাহু ধরলে, সেই বিন্দু থেকে আঁকা কর্ণই লব্ধি।}
\sqitem{মান:} $R=\sqrt{P^2+Q^2+2PQ\cos\theta}$
\sqitem{দিক:} $\tan\alpha=\dfrac{Q\sin\theta}{P+Q\cos\theta}$
\begin{roundbox}
\textbullet\ কোণ যত বাড়ে, লব্ধি তত কমে।\par
\textbullet\ সমান মানের দুটি ভেক্টরের লব্ধি সবসময় তাদের \textbf{কোণের সমদ্বিখণ্ডক} বরাবর।
\end{roundbox}
\end{typebox}

\begin{tipbox}[455]
$\alpha$ সবসময় যে ভেক্টরের সাথে মাপছো, সূত্রে সেটাকেই $P$ ধরো — না হলে দিক উল্টো আসবে। পরীক্ষায় "কার সাথে" লিখতে ভুলবে না।
\end{tipbox}

\cq{১}[ভিডিও ম্যাথ]{৬ একক ও ৮ একক মানের দুটি ভেক্টর $60^\circ$ কোণে ক্রিয়া করছে। লব্ধির মান ও দিক নির্ণয় করো।}[520]
\solhead
\workgiven[0.62]{%
$R=\sqrt{P^2+Q^2+2PQ\cos\theta}$\\
$\;\;=\sqrt{6^2+8^2+2\cdot6\cdot8\cos60^\circ}$\\
$\;\;=\sqrt{36+64+48}=\sqrt{148}$\\
$\;\;=12.17$ একক \Ans\par\vspace{2pt}
$\tan\alpha=\dfrac{8\sin60^\circ}{6+8\cos60^\circ}=\dfrac{6.928}{10}$\\
$\alpha=34.71^\circ$ \Ans}
{\given{$P=6$ একক\\ $Q=8$ একক\\ $\theta=60^\circ$}}

\cq{২}[ভিডিও ম্যাথ]{সমান মানের দুটি বল $F$ এর লব্ধির মানও $F$ হলে বল দুটির মধ্যবর্তী কোণ কত?}[760]
\solhead
\workgiven[0.62]{%
$F^2=F^2+F^2+2F^2\cos\theta$\\
$\Rightarrow\cos\theta=-\dfrac12$\\
$\therefore\ \theta=120^\circ$ \Ans}
{\given{$P=Q=F$\\ $R=F$}}

% ---------------------------------------------------------------- TYPE 02
\begin{typebox}{02}{আপেক্ষিক বেগ: নদী–নৌকা}
\withqr{1120}{\textbf{কনসেপ্ট:} নৌকার লব্ধি বেগ = নৌকার নিজস্ব বেগ $\vec v$ + স্রোতের বেগ $\vec u$। কোন দিকে চালাবে তা নির্ভর করে কী চাই তার উপর — \textbf{কম পথ} নাকি \textbf{কম সময়}।}
\begin{center}
\begin{tikzpicture}[scale=.82]
  \fill[accent!8] (0,0) rectangle (5.4,2.4);
  \draw[line width=1pt] (0,0)--(5.4,0) (0,2.4)--(5.4,2.4);
  \foreach \y in {.6,1.2,1.8}{\draw[vecthin,accent!60] (3.7,\y)--(4.9,\y);}
  \node[lbl,accent] at (4.3,2.1) {স্রোত $\vec u$};
  \coordinate (A) at (1.6,0);
  \draw[vec] (A)--++(-1.2,2.4) node[above,lbl]{$\vec v$};
  \draw[vec,accent] (A)--++(0,2.4) node[above,lbl]{$\vec R$};
  \draw[dim] (5.65,0)--(5.65,2.4) node[midway,right,lbl]{$d$};
  \draw (1.6,.7) arc (90:116.6:.7) node[midway,above,lbl]{$\beta$};
\end{tikzpicture}
\end{center}
\sqitem{সর্বনিম্ন পথ (সোজা পার):} $\sin\beta=\dfrac uv,\quad t=\dfrac{d}{\sqrt{v^2-u^2}}$
\sqitem{সর্বনিম্ন সময় (লম্বভাবে চালাও):} $t_{\min}=\dfrac dv,\quad x=u\,t_{\min}$
\end{typebox}

\cq{৩}[ভিডিও সৃজনশীল]{$1\,\text{km}$ প্রশস্ত নদীতে স্রোতের বেগ $3\,\text{km\,h}^{-1}$। স্থির পানিতে নৌকার বেগ $5\,\text{km\,h}^{-1}$।}[1310]
\qpart{গ}{সবচেয়ে কম পথে পার হতে কত সময় লাগবে?}
\qpart{ঘ}{সবচেয়ে কম সময়ে পার হলে নৌকা কতদূর ভাটিতে সরে যাবে? গাণিতিক বিশ্লেষণ দাও।}
\solhead
\anslabel{গ}
\workgiven[0.62]{%
$R=\sqrt{v^2-u^2}=\sqrt{25-9}$\\
$\;\;=4\,\text{km\,h}^{-1}$\\
$t=\dfrac dR=\dfrac14\,\text{h}=15$ মিনিট \Ans}
{\given{$d=1\,\text{km}$\\ $u=3\,\text{km\,h}^{-1}$\\ $v=5\,\text{km\,h}^{-1}$}}
\anslabel{ঘ}
নৌকাকে স্রোতের লম্বভাবে চালাতে হবে।\par
$t_{\min}=\dfrac dv=\dfrac15\,\text{h}=12$ মিনিট\par
$x=u\,t_{\min}=3\times\dfrac15=0.6\,\text{km}$\par
$\therefore$ নৌকা $600\,\text{m}$ ভাটিতে সরে যাবে। \Ans

% ---------------------------------------------------------------- TYPE 03
\begin{typebox}{03}{ডট ও ক্রস গুণন}
\withqr{1865}{\textbf{কনসেপ্ট:} ডট গুণনে ফল একটি \textbf{স্কেলার} (সংখ্যা), ক্রস গুণনে ফল একটি \textbf{ভেক্টর} যার দিক দুই ভেক্টরের তলের লম্ব বরাবর (ডানহাতি স্ক্রু নিয়ম)।}
\sqitem{কোণ:} $\cos\theta=\dfrac{A_xB_x+A_yB_y+A_zB_z}{AB}$
\sqitem{ক্রস গুণন নির্ণায়ক:}
$\vec A\times\vec B=\begin{vmatrix}\hat i&\hat j&\hat k\\A_x&A_y&A_z\\B_x&B_y&B_z\end{vmatrix}$
\end{typebox}

\begin{tipbox}[2010]
ডট গুণনে ঋণাত্মক উত্তর আসলে ঘাবড়াবে না — মানে কোণটা $90^\circ$ এর বেশি (স্থূলকোণ)।
\end{tipbox}

\cq{৪}[ভিডিও ম্যাথ]{$\vec A=2\hat i+3\hat j-\hat k$ এবং $\vec B=\hat i-2\hat j+2\hat k$ হলে এদের মধ্যবর্তী কোণ ও $|\vec A\times\vec B|$ নির্ণয় করো।}[2140]
\solhead
$\vec A\cdot\vec B=(2)(1)+(3)(-2)+(-1)(2)=-6$\par
$A=\sqrt{4+9+1}=\sqrt{14},\quad B=\sqrt{1+4+4}=3$\par
$\cos\theta=\dfrac{-6}{3\sqrt{14}}=-0.5345$\par
$\therefore\ \theta=122.31^\circ$ \Ans\par\vspace{2pt}
$\vec A\times\vec B=\begin{vmatrix}\hat i&\hat j&\hat k\\2&3&-1\\1&-2&2\end{vmatrix}=4\hat i-5\hat j-7\hat k$\par
$|\vec A\times\vec B|=\sqrt{16+25+49}=\sqrt{90}=9.49$ একক \Ans
\end{multicols}
\end{document}
